\chapter*{Glosario}
\addcontentsline{toc}{chapter}{Glosario}
\indicetoprule

Este glosario reúne los términos de mercado y decisión empleados de forma recurrente en la memoria.
\begin{description}
  \item[\glostarget{activo-digital}{Activo digital}] Elemento asociado a una cuenta de usuario o a un servicio digital que puede utilizarse, coleccionarse o intercambiarse.
  \item[\glostarget{accion-conjunta}{Acción conjunta}] Conjunto formado por las acciones que todos los agentes de un entorno multiagente toman en un mismo instante.
  \item[\glostarget{agente}{Agente}] Programa que toma decisiones a partir de la información que recibe.
  \item[\glostarget{aprendizaje-refuerzo}{Aprendizaje por refuerzo}] Método de aprendizaje en el que un agente observa un entorno, selecciona acciones y recibe recompensas que permiten evaluar sus consecuencias.
  \item[\glostarget{aprendizaje-refuerzo-multiagente}{Aprendizaje por refuerzo multiagente}] Extensión del aprendizaje por refuerzo en la que dos o más agentes interactúan con un mismo entorno y pueden cooperar, competir o decidir de forma independiente.
  \item[\glostarget{beneficio-neto}{Beneficio neto}] Resultado económico de una operación después de descontar el precio de compra, las comisiones y los costes aplicables.
  \item[\glostarget{capital-bloqueado}{Capital bloqueado}] Parte del saldo que permanece comprometida en una compra o que no puede reutilizarse mientras el activo no pueda venderse o transferirse.
  \item[\glostarget{cartera}{Cartera}] Conjunto formado por el saldo disponible y las posiciones abiertas, es decir, los activos digitales adquiridos que aún no se han vendido.
  \item[\glostarget{cookie-sesion}{Cookie de sesión}] Dato local generado por una plataforma web que permite reconocer una sesión iniciada sin volver a introducir las credenciales en cada petición.
  \item[\glostarget{demanda}{Demanda}] Número de personas usuarias que quieren adquirir un activo digital.
  \item[\glostarget{disponibilidad}{Disponibilidad}] Número de unidades de un activo digital que están a la venta en una plataforma.
  \item[\glostarget{economia-videojuego}{Economía de videojuego}] Conjunto de reglas con las que un videojuego permite obtener, poseer, utilizar e intercambiar activos digitales.
  \item[\glostarget{economia-persistente}{Economía persistente}] Economía de videojuego en la que los activos digitales permanecen asociados a la cuenta de la persona usuaria después de terminar una partida.
  \item[\glostarget{entorno}{Entorno}] Representación de la situación en la que un agente toma decisiones, formada por la información disponible y las reglas que determinan las consecuencias de sus acciones.
  \item[\glostarget{episodio}{Episodio}] Ejecución completa de un entorno de aprendizaje por refuerzo, desde su estado inicial hasta la condición que marca su finalización.
  \item[\glostarget{escasez}{Escasez}] Número de unidades existentes de un activo digital en todo el mundo.
  \item[\glostarget{estado-global}{Estado global}] Información completa que el entorno utiliza en un instante para calcular observaciones, transiciones y recompensas.
  \item[\glostarget{exposicion}{Exposición}] Parte de la cartera comprometida con un activo, una plataforma o una categoría de activos.
  \item[\glostarget{funcion-transicion}{Función de transición}] Regla del entorno que determina cómo cambia el estado después de ejecutar una acción.
  \item[\glostarget{horizonte-temporal}{Horizonte temporal}] Periodo previsto entre la compra de un activo y su posible venta.
  \item[\glostarget{mercado-cerrado}{Mercado cerrado}] Entorno en el que los activos solo se utilizan dentro del videojuego y no se intercambian entre personas usuarias.
  \item[\glostarget{mercado-centralizado}{Mercado centralizado}] Mercado en el que una empresa o plataforma administra las cuentas, los activos digitales o las reglas de intercambio.
  \item[\glostarget{mercado-descentralizado}{Mercado descentralizado}] Entorno en el que la gestión del activo y las reglas de intercambio no dependen de una plataforma central.
  \item[\glostarget{mercado-digital}{Mercado digital}] Entorno en línea en el que las personas usuarias intercambian bienes o servicios mediante plataformas.
  \item[\glostarget{mercado-integrado}{Mercado integrado}] Mercado administrado por el propio ecosistema del videojuego o de distribución, que controla los activos digitales asociados a las cuentas y las reglas de intercambio.
  \item[\glostarget{mercado-secundario}{Mercado secundario}] Mercado de reventa entre personas usuarias.
  \item[\glostarget{mypy}{Mypy}] Herramienta de comprobación estática de tipos para Python que revisa si el uso de variables y funciones es coherente con sus anotaciones.
  \item[\glostarget{oferta}{Oferta}] Publicación mediante la cual una persona pone a la venta un activo digital identificado por sus características relevantes e indica su precio.
  \item[\glostarget{operacion}{Operación}] Decisión sobre un activo digital. Puede consistir en comprar, vender, mantener una posición abierta o descartar una compra o una venta.
  \item[\glostarget{operacion-compraventa}{Operación de compraventa}] Análisis conjunto de una compra de un activo digital y su venta posterior.
  \item[\glostarget{oportunidad}{Oportunidad}] Posible operación de compraventa que merece revisarse porque sus datos iniciales indican que podría generar un beneficio neto.
  \item[\glostarget{observacion-local}{Observación local}] Parte del estado global que recibe un agente concreto para tomar su decisión.
  \item[\glostarget{orden-compra}{Orden de compra}] Publicación que indica el precio máximo que una persona está dispuesta a pagar por un activo digital concreto.
  \item[\glostarget{plataforma}{Plataforma}] Servicio digital donde se administran las reglas del intercambio y la información necesaria para participar en él.
  \item[\glostarget{plataforma-externa}{Plataforma externa}] Plataforma distinta del servicio que administra originalmente los activos digitales vinculados a la cuenta y que publica ofertas o facilita operaciones sobre activos transferibles.
  \item[\glostarget{politica}{Política}] Regla o modelo que indica qué acción selecciona un agente a partir de la información que observa.
  \item[\glostarget{posicion-abierta}{Posición abierta}] Activo digital ya adquirido que permanece en la cartera hasta su venta. Mientras la posición está abierta, el importe empleado en la compra no puede destinarse a otra operación.
  \item[\glostarget{pytest}{Pytest}] Herramienta de pruebas para Python que localiza funciones de prueba, las ejecuta y comunica si el comportamiento observado coincide con el esperado.
  \item[\glostarget{rareza}{Rareza}] Dificultad para obtener un activo digital.
  \item[\glostarget{recompensa}{Recompensa}] Señal numérica que evalúa la consecuencia de una acción en aprendizaje por refuerzo y guía el aprendizaje de una política.
  \item[\glostarget{rentabilidad}{Rentabilidad}] Relación entre el beneficio neto de una operación y el capital empleado para realizarla.
  \item[\glostarget{restriccion-riesgo}{Restricción de riesgo}] Regla que limita una operación para evitar que la cartera concentre demasiado capital, pierda liquidez o incumpla condiciones operativas.
  \item[\glostarget{ruff}{Ruff}] Herramienta para Python que revisa reglas de estilo, errores estáticos y formato de código.
  \item[\glostarget{senal-supervisada}{Señal supervisada}] Valor producido por un modelo entrenado con ejemplos históricos y utilizado como información adicional antes de tomar una decisión.
  \item[\glostarget{scraping}{Scraping}] Proceso de consulta automática de una fuente web para extraer datos y convertirlos en registros estructurados.
  \item[\glostarget{spread}{Spread}] Diferencia entre el mejor precio de venta y la mejor orden de compra disponibles para un mismo activo en un momento dado.
  \item[\glostarget{trade-hold}{Trade hold}] Bloqueo temporal que impide vender o transferir un activo digital recién adquirido.
  \item[\glostarget{volumen-intercambios}{Volumen de intercambios}] Cantidad de unidades de un activo digital vendidas durante un periodo, como un día, una semana o un mes.
\end{description}
\indicebottomrule
